%==================================================================

\subsection{Ядро и образ}
\leavevmode\par
\begin{defn}[ядро и образ]
Пусть $f\colon\Shv{F}\to\Shv{G}$ --- гомоморфизм пучков.
\emph{Ядром} $\ker f$ называется пучок, определяемый поэлементно:
\begin{equation}\label{eq:ker}
 (\ker f)(U)=\ker\bigl(\Shv{F}(U)\xrightarrow{f_U}\Shv{G}(U)\bigr),
\end{equation}
ограничения --- обычные. \emph{Образ} определяется по формуле
\begin{equation}\label{eq:im}
 \operatorname{Im}f=\bigl(\operatorname{Im}f\bigr)^{\#},
\end{equation}
где $\operatorname{Im}f$ справа --- предпучок с группами
$\operatorname{Im}\bigl(\Shv{F}(U)\xrightarrow{f_U}\Shv{G}(U)\bigr)$,
а $({-})^{\#}$ --- прокачивание.
\end{defn}

Почему образ \emph{нельзя} оставить поэлементным: множество
$\operatorname{Im}(f_U)$, вообще говоря, не является пучком. Пример:
пусть $\Shv{F}$ --- постоянный предпучок $\mathbb{Z}$, а
$\Shv{G}$ --- пучок, в котором локально ненулевые вещи склеиваются;
образ поэлементно будет <<предпучком, который надо досклеить>>.
Прокачивание~\eqref{eq:im} --- это и есть правильный образ: он
удовлетворяет (F2) по построению.

\begin{rem}[ядро не требует прокачивания]
Ядро --- обратная операция: точность поэлементна сохраняется, потому
что $\Shv{F}\mapsto\Shv{F}(U)$ --- точный функтор (см.~Годемана,
п.~1.9). Поэтому $\ker f$ поэлементно --- уже пучок.
\end{rem}

\subsection{Точные последовательности}
\leavevmode\par
\begin{defn}
Последовательность пучков
$\Shv{F}\xrightarrow{\ f\ }\Shv{G}\xrightarrow{\ g\ }\Shv{H}$
называется \emph{точной}, если $\ker g=\operatorname{Im}f$.
\end{defn}

Диаграмма~\ref{dgm:canon} показывает каноническое разложение
гомоморфизма, из которого и состоит определение.

\begin{diagram}{Каноническое разложение гомоморфизма пучков и
 точность: $\operatorname{Im}f=\ker g$.\label{dgm:canon}}
\begin{tikzcd}[cd style,row sep=large]
0 \ar[r] & \ker f \ar[r,hook] & \Shv{F} \ar[r,"f"'] \ar[d,two heads] &
\Shv{G} \ar[r,"g"] & \Shv{H} \\
 & & \operatorname{Im} f \ar[r,equal] & \ker g \ar[r] \ar[u,hook] & 0
\end{tikzcd}
\end{diagram}

Пояснение к диаграмме~\ref{dgm:canon}. Две вертикали и стрелка
равенства и есть каноническое разложение: $\Shv{F}$ сюръективно
спускается на свой образ $\operatorname{Im}f$ (вниз, <<две
головы>>), образ отождествляется с $\ker g$, а $\ker g$
включается обратно в $\Shv{G}$ (вверх, крючок). Таким образом
$f$ читается как <<сначала спуститься на образ, потом включить>>:
$\Shv{F}\twoheadrightarrow\operatorname{Im}f=\ker g\hookrightarrow
\Shv{G}$. Слева $\ker f\hookrightarrow\Shv{F}$ --- участок, на
котором $f$ равен нулю; полнота (инъективность слева,
сюръективность справа) как раз закодирована нулями по краям.

\subsection{Как проверять точность}
Практическое правило: точность последовательности пучков проверяется
на ростках $\Shv{F}_x\to\Shv{G}_x\to\Shv{H}_x$ для всех $x\in X$.
Эквивалентно, каждое сечение из ядра второго отображения должно
\emph{локально} на покрытии быть образом сечения первого пучка.
Не следует требовать поточечной сюръективности отображений сечений над
каждым открытым множеством: эпиморфизм пучков вообще не обязан быть
сюръективен на глобальных сечениях.

Частый пример: $0\to\Shv{F}\to\Shv{G}$ точна тогда и только тогда,
когда $\Shv{F}\to\Shv{G}$ --- инъекция \emph{в каждом слое}.

\begin{bookbox}
\booksrc{Годеман, гл.~I, п.~1.8, с.~26---29 --- ядро, коядро, образ и
 кообраз в аддитивной категории, определение точной последовательности;
 гл.~II, \S\,2.4---2.5, с.~153---154 --- каноническое разложение
 гомоморфизма пучков и точные последовательности: точность
 последовательности $\mathcal{O}$-модулей эквивалентна точности на
 каждой точке $x$, т.\,е. $\ker(g)(x)=\operatorname{Im}(f)(x)$.
\par\smallskip
 Манин, \S\,2.1.4, с.~110 --- аксиома пучка как точность
 последовательности
 $0\to\Shv{P}(U)\to\prod_i\Shv{P}(U_i)\rightrightarrows
 \prod_{i,j}\Shv{P}(U_i\cap U_j)$;
 \S\,2.9, с.~141 --- обзор предпучков и пучков модулей.
}
\end{bookbox}

%==================================================================
